\section{Introduction}

LLM evaluation frameworks measure trustworthiness dimensions independently---truthfulness, robustness, fairness, safety, privacy---yet deployment failures often span multiple dimensions simultaneously. A model failing a factual question may also fail robustness tests on the same instance, but current benchmarks provide no visibility into these cross-dimensional coupling patterns. This gap has concrete consequences: a medical diagnosis LLM might achieve 92\% on truthfulness benchmarks and 89\% on fairness benchmarks when evaluated separately, yet fail on both dimensions for the same minority-group patients---a coupling pattern invisible to independent evaluation.

The stakes are high. Independent dimension evaluation creates blind spots for compound failures---instances where models fail on multiple trustworthiness criteria at once. Understanding coupling patterns is critical for model selection in safety-critical deployments. Without coupling awareness, deployment teams select models based on aggregate scores that mask correlated vulnerabilities, leading to systematic failures in production.

Two hypotheses dominate the literature. The \textbf{broad independence hypothesis} (implicit in TrustLLM and MMTrustEval frameworks) assumes zero coupling---dimensions reflect orthogonal model capabilities, so evaluating them independently is sufficient. The \textbf{broad coupling hypothesis} (suggested by Li \& Li 2024 documenting robustness-fairness triangular trade-offs) predicts pervasive co-occurrence from shared architectural bottlenecks, where stress causes cascading failures across many dimension pairs. Neither hypothesis has been empirically tested across all dimension pairs.

Multi-dimensional trustworthiness evaluation exists. Frameworks like TrustLLM and MMTrustEval cover 5--8 dimensions, providing per-dimension scores for model comparison. Yet this independent evaluation lacks co-occurrence statistics showing which instances fail on multiple dimensions simultaneously. Shared vulnerability mechanisms (e.g., calibration failures affecting both truthfulness and robustness) may create systematic coupling patterns that current benchmarks cannot detect.

The deeper problem: independent evaluation may miss model-specific vulnerability architectures. Coupling patterns themselves may be a trustworthiness property---model-specific fingerprints revealing underlying vulnerability structure. If different models exhibit distinct coupling profiles (GPT-4 couples truthfulness-robustness, Claude couples fairness-safety), these patterns could guide deployment decisions for specific use cases. But without instance-level binary labels across all dimension pairs, we cannot answer: (1) Are dimensions fundamentally independent or coupled? (2) Do different models have distinct vulnerability profiles? (3) Can coupling guide model selection?

We address this gap by demonstrating a methodology framework for measuring coupling breadth and model-specificity in LLM trustworthiness dimensions. Our key finding from synthetic validation: coupling is \textbf{sparse}---limited to 2 dominant dimension pairs (truthfulness-robustness phi 0.36--0.40, fairness-safety phi 0.33--0.40), not pervasive across all 10 possible pairs. This sparsity persists when controlling for instance difficulty (partial phi 0.36--0.56, retention 86--161\%), indicating genuine shared vulnerabilities rather than spurious artifacts of hard instances failing everywhere. Models show distinct coupling profiles (GPT-4: truthfulness$\rightarrow$robustness chain, Claude-3: fairness$\rightarrow$safety$\rightarrow$privacy cluster, Llama-3: minimal coupling) in qualitative patterns, though statistical confirmation requires larger samples than our Phase 4 proof-of-concept (n=100/dimension).

\textbf{Critical limitation:} All experiments used synthetic coupling data designed to emulate benchmark structure. This validates measurement methodology but defers real-world coupling characterization to Phase 5 production benchmark evaluation. Coupling patterns (phi values, sparsity, profiles) shown here demonstrate detectability IF coupling exists, not that these specific magnitudes exist in real LLMs.

Our contributions build on this sparse coupling methodology demonstration:

\begin{enumerate}
\item \textbf{Coupling breadth measurement framework} via phi coefficient analysis across 5 dimensions $\times$ 3 models (h-e1: 6 significant pairs detectable at phi 0.33--0.40, $p < 1\times10^{-13}$), establishing coupling as a phenomenon measurable from behavioral outputs when present.

\item \textbf{Difficulty-independence validation framework} via partial correlation controlling for instance difficulty (h-m1: partial phi 0.36--0.56, retention 86--161\%), proving methodology can isolate shared vulnerability mechanisms from confounding by hard instances.

\item \textbf{Sparsity detection capability} via multi-pair analysis with Bonferroni correction (h-m2: 0 models meet $\geq$3 pairs threshold in synthetic data), demonstrating methodology can distinguish sparse from broad coupling patterns.

\item \textbf{Model-specific profile observation} via Mantel test comparing coupling matrices (h-c1: qualitative patterns observable but statistical confirmation underpowered at n=100; requires $n \geq 500$), providing suggestive evidence that methodology could detect architectural fingerprints given adequate samples.
\end{enumerate}

Pending real-world validation, these patterns would reframe multi-dimensional trustworthiness evaluation. Rather than assuming universal independence (coupling = 0) or universal coupling (all pairs correlated), our methodology reveals when trustworthiness dimensions exhibit \textbf{sparse coupling}---architecturally independent by default, with coupling emerging only where specific mechanisms overlap (calibration for truthfulness-robustness, value alignment for fairness-safety). If validated, benchmarks should report coupling statistics alongside per-dimension scores, and model selection for deployments requiring multiple trustworthiness guarantees can exploit known coupling patterns.

Our work builds on multi-dimensional evaluation frameworks and behavioral detection methods, but adds a coupling measurement layer absent from prior work. The next section positions our approach relative to existing trustworthiness evaluation and vulnerability analysis literature.
